%======================================================================
\section{Positivity on the Fourier side}\label{sec:fourier}
%======================================================================

On $(1,\infty)$ the $\Gamma$-only transform is $\sin^2(\pi x)\sqrt x\sum_na_nK_0(4\pi\sqrt{nx})$ (Theorem~\ref{thm:lift}(c)), so the sign of
$\Gh_{\rm raw}$ beyond $x=1$ is decided by the signs of the coefficients $a_n=n\Ss_n+\frac14\delta_{n1}$ of \eqref{eq:an}. This section proves that
every $a_n$ is positive, with explicit two-sided bounds, using only the trivial bound for Kloosterman sums. It follows that $\Xi^2H_{\rm raw}$ is
positive definite. This gives the sign of the normalising constant without computation (Section~\ref{ssec:C}), and, through the moments of the
transform of $\Xi^2H_{\rm raw}$, the positivity of $H_{\rm raw}$ for $|t|\le10.355$ (Section~\ref{ssec:moments}).

\subsection{The coefficients are positive}

Put $z_n:=4\pi\sqrt n$, $\ee(w):=e^w/\sqrt{2\pi w}$ and $\abar_n:=n\cdot\frac2{z_n}\ee(z_n)=\frac{n^{1/4}e^{4\pi\sqrt n}}{4\sqrt2\,\pi^2}$. With
$\kappa_1:=1+e/3$, $\kappa_2:=2e/9$ and $\kappa_3:=\kappa_1\psi(3)+\kappa_2=2.36297\ldots$, define for $z\ge4\pi$
\begin{gather*}
\lambda_0(z):=\Big(1-\frac1z\Big)\Big(1-\frac3{8z}-\frac{15}{32z^2}\Big)-\frac{\pi^3}{16z^2},\\
\varrho(z):=\frac{(\pi+2)z}{8\ee(z)}+\sqrt2e^{-z/2}+\frac{(\pi+2)z^2}{8\ee(z)}+\frac{\sqrt3z^2}4e^{-2z/3}+\Big[\kappa_3+\frac z2(\kappa_1+\kappa_3)\Big]\frac z{2\ee(z)},\\
\lambda_-(z):=\lambda_0(z)-\frac{(\pi+1)z}{4\ee(z)}-\varrho(z),\qquad \lambda_+(z):=1-\frac1z+\Big(\frac{\pi+1}2+\frac2{z^2}\Big)\frac z{2\ee(z)}+\varrho(z).
\end{gather*}

\begin{theorem}[Positivity of the coefficients]\label{thm:an-positive}
For every $n\ge1$,
\[
\frac2{z_n}\ee(z_n)\,\lambda_-(z_n)\le\Ss_n\le\frac2{z_n}\ee(z_n)\,\lambda_+(z_n),
\]
and $\lambda_-$ is increasing on $[4\pi,\infty)$ with $\lambda_-(4\pi)=0.85010\ldots$. Hence $\Ss_n>0$ and $a_n>0$ for every $n\ge1$.
\end{theorem}

\begin{theorem}[Two-sided bounds]\label{thm:an-bounds}
For every $n\ge1$, $(1-2/z_n)\,\abar_n\le a_n\le\abar_n$.
\end{theorem}

The proofs use only the trivial bound $|S(\pm1,n;c)|\le\varphi(c)$ and involve no check over finitely many $n$. Their only numerical inputs are
elementary evaluations at the single point $z=4\pi$: $\lambda_0(4\pi)=0.877952$, $(\pi+1)z/(4\ee(z))=0.000403$, and the five
summands of $\varrho(4\pi)$ are $0.000250$, $0.002641$, $0.003145$, $0.015725$, $0.005683$, so $\lambda_-(4\pi)=0.850105$. We evaluated all of them
in ball arithmetic as a check (Appendix~\ref{app:cert}); a pocket calculator suffices.

\begin{proof}[Proof of Theorem~\ref{thm:an-positive}]
Fix $n$ and write $z=z_n\ge4\pi$, $M(w):=|\Jdot(w)|+|\Idot(w)|$, and $\Ss_n=T_1+R_n$ with $T_1:=\Jdot(z)-\Idot(z)$ the term $c=1$. By the trivial
bound, $|R_n|\le\sum_{c\ge2}(\varphi(c)/c)M(z/c)$. We split the sum, using the Bessel bounds of Lemmas~\ref{lem:Jdot-uniform}--\ref{lem:Idot} in
Appendix~\ref{app:bessel}.
\begin{enumerate}[label=(\roman*)]
\item $c=2$: $\varphi(2)/2=\frac12$ and $w=z/2\ge2\pi$, so $M(z/2)\le\frac{\pi+2}2+\frac4z\ee(z/2)$; since $\ee(z/2)/\ee(z)=\sqrt2e^{-z/2}$, the contribution
divided by $\frac2z\ee(z)$ is at most $\frac{(\pi+2)z}{8\ee(z)}+\sqrt2e^{-z/2}$.
\item $3\le c\le z/2$: $w=z/c\ge2$ and $\varphi(c)/c\le1$, so $M(z/c)\le\frac{\pi+2}2+\frac{2c}z\ee(z/c)\le\frac{\pi+2}2+\ee(z/3)$, as $\ee$ is increasing on
$[\frac12,\infty)$. There are at most $z/2$ such $c$, and $\ee(z/3)/\ee(z)=\sqrt3e^{-2z/3}$; the contribution divided by $\frac2z\ee(z)$ is at most
$\frac{(\pi+2)z^2}{8\ee(z)}+\frac{\sqrt3z^2}4e^{-2z/3}$.
\item $c>z/2$: $w=z/c<2$, and by Lemma~\ref{lem:small-w}(c), $(\varphi(c)/c)M(z/c)\le f(c)$ with
$f(y):=\frac{z^2}4y^{-2}[\kappa_1\log(2y/z)+\kappa_3]$, which is positive and decreasing on $[z/2,\infty)$ (as $\kappa_1<2\kappa_3$). So
$\sum_{c>z/2}f(c)\le f(z/2)+\int_{z/2}^\infty f=\kappa_3+\frac z2(\kappa_1+\kappa_3)$, and divided by $\frac2z\ee(z)$ this is the last summand of $\varrho$.
\end{enumerate}
Hence $|R_n|\le\frac2z\ee(z)\varrho(z)$. For the term $c=1$, Lemmas~\ref{lem:Jdot-uniform} and~\ref{lem:Idot}(c) give
$T_1\ge\frac2z\ee(z)\lambda_0(z)-\frac{\pi+1}2$ and $T_1\le\frac{\pi+1}2+\frac2{z^2}+\frac2z\ee(z)(1-\frac1z)$; this gives both inequalities.
\emph{Monotonicity.} $\lambda_0$ is increasing on $[1,\infty)$: it is a product of two non-negative increasing factors minus a decreasing term.
Each of $z^k/\ee(z)\propto z^{k+1/2}e^{-z}$ ($k\le3$), $e^{-z/2}$ and $z^2e^{-2z/3}$ is decreasing on $[4\pi,\infty)$, so every subtracted term in
$\lambda_-$ decreases there, and $\lambda_-$ increases. Finally $\lambda_-(4\pi)=0.850105>0$, so $\Ss_n>0$, and $a_n=n\Ss_n+\frac14\delta_{n1}>0$.
\end{proof}

\begin{proof}[Proof of Theorem~\ref{thm:an-bounds}]
Upper bound: $z(\lambda_+(z)-1+\frac1z)=(\frac{\pi+1}2+\frac2{z^2})\frac{z^2}{2\ee(z)}+z\varrho(z)$ is decreasing on $[4\pi,\infty)$, by the same monotonicity
facts, and equals $0.34996$ at $z=4\pi$; so $\lambda_+(z)\le1-0.65/z$ and $a_n\le\abar_n(1-0.65/z_n)+\frac14\delta_{n1}\le\abar_n$, because
$0.65\,\abar_1/(4\pi)=265.6>\frac14$. Lower bound: $a_n\ge n\Ss_n\ge\abar_n\lambda_-(z_n)$, and
$z(\lambda_-(z)-1+\frac2z)=z(\lambda_0(z)-1+\frac2z)-z[\frac{(\pi+1)z}{4\ee(z)}+\varrho(z)]$, where
$z(\lambda_0(z)-1+\frac2z)=0.625-(\frac3{32}+\frac{\pi^3}{16})z^{-1}+\frac{15}{32}z^{-2}$ is increasing (value $0.46630$ at $4\pi$) and the bracketed term
is decreasing (value $0.34993$ at $4\pi$). So $\lambda_-(z)\ge1-2/z$ on $[4\pi,\infty)$.
\end{proof}

The bounds are sharp in the main term: the term $c=1$ alone is $nT_1=\frac{\sqrt n}{2\pi}I_1(z_n)-\frac{I_0(z_n)}{8\pi^2}+n[K_2(z_n)+\Jdot(z_n)]$ by
Lemma~\ref{lem:bessel-basic}(c), and the certified values of Appendix~\ref{app:cert} show $\sum_{c\ge2}|\text{term}_c|\le3.6\cdot10^{-3}\,T_1$ for
$n\le1500$, so Theorem~\ref{thm:an-positive} loses only a factor of about $8$ in $\varrho$.

\emph{The trivial-bound box.} The proof of Theorem~\ref{thm:an-positive} gives more. With $T_1(z)=\Jdot(z)-\Idot(z)$ and $R_n$ as in that proof,
$nT_1(z_n)$ is the term $c=1$ of $n\Ss_n$; put
\begin{equation}\label{eq:box}
c_n:=nT_1(z_n)+\tfrac14\delta_{n1},\qquad w_n:=\abar_n\varrho(z_n).
\end{equation}
Then $|a_n-c_n|=n|R_n|\le w_n$ for every $n$. We call the set $\TB$ of real sequences $(b_n)_{n\ge1}$ with $|b_n-c_n|\le w_n$ for every $n$ the
\emph{trivial-bound box}; that $(a_n)\in\TB$ uses only $S(\pm1,n;1)=1$ and $|S(\pm1,n;c)|\le\varphi(c)$. By the two proofs above, every
$(b_n)\in\TB$ satisfies $(1-2/z_n)\abar_n\le\abar_n\lambda_-(z_n)\le b_n\le\abar_n\lambda_+(z_n)+\frac14\delta_{n1}\le\abar_n$. The relative
half-width $w_n/c_n$ is $3.09\cdot10^{-2}$, $1.40\cdot10^{-3}$ and $1.41\cdot10^{-4}$ for $n=1,2,3$. The two certificates on which the positivity
of $H$ rests, Theorem~\ref{thm:moments} and Proposition~\ref{prop:L-numbers-box}, use the coefficients only through $(a_n)\in\TB$: each bound they
certify holds for every sequence in $\TB$ in place of $(a_n)$.

\begin{remark}[The Weil bound is not needed]
The Weil--Estermann bound $|S(m,n;c)|\le d(c)(m,n,c)^{1/2}c^{1/2}$ \cite{Wei48,Est61} (see also \cite[Chapter~11]{IK04}) would improve the constants, but
the trivial bound already suffices for every $n$. The Weil bound enters the paper only through the tails of the certified coefficient
enclosures (Appendix~\ref{app:cert-coeffs}). In the proofs of the main results these are used only for the value of $C$
(Proposition~\ref{prop:C}); they also enter the independent checks of Section~\ref{sec:physical}.
\end{remark}

\subsection{The normalising constant}\label{ssec:C}

The function $H_{\rm raw}$ of Definition~\ref{def:Hraw} has $\Gamma$-only transform $\Gh_{\rm raw}$, whose coefficients are the $a_n$. The object of
Theorem~\ref{thm:main-object} is normalised by $H(0)=1$, i.e.\ $H:=H_{\rm raw}/H_{\rm raw}(0)$, provided $H_{\rm raw}(0)\neq0$. The sign of
$H_{\rm raw}(0)$ follows from the positivity of the $a_n$ without computation (Proposition~\ref{prop:posdef} below); its value is
computer-assisted.

\begin{proposition}[Computer-assisted]\label{prop:C}
$H_{\rm raw}(0)=0.0011355097904\pm1.86\cdot10^{-9}$. Hence $C:=1/H_{\rm raw}(0)\in[880.66027,\,880.66316]$, and
$\Gh_{\rm raw}(0)\cdot C=C/(32\pi^2)=2.788428\pm5\cdot10^{-6}$.
\end{proposition}

The certificate (Appendix~\ref{app:cert-C}) evaluates the Mellin transform of $\Gh_{\rm raw}/(2\pi\sqrt x)$ at $s=\frac12$ through a route that is
independent of the physical-side formula of Section~\ref{sec:physical}: the density \eqref{eq:D-cusp0} on $(0,1]$, an explicit cusp-$\infty$
expansion of $D$ on $[1,\infty)$, and closed-form Mellin kernels. The two routes share only the coefficient enclosures. Through the
physical-side formula, $H_{\rm raw}(0)$ is very sensitive to $a_1$ (Remark~\ref{rem:conditioning}); the route used here is far less so, which
is why its radius is small. With
Theorem~\ref{thm:an-bounds}, Proposition~\ref{prop:C} gives the normalised bounds of Theorem~\ref{thm:main-object}:
$C/(4\sqrt2\pi^2)\in[15.77370,15.77376]$.

\begin{proposition}[Positive definiteness]\label{prop:posdef}
Put $F_{\rm raw}:=\Xi^2H_{\rm raw}$. For every $\xi\ge0$,
\[
\widehat{F_{\rm raw}}(\xi)=\sum_{m\ge1}d(m)\,m^{-1/2}\,\Gh_{\rm raw}(\xi+\xin m)\ge0,
\]
where the $m$-th argument is the point $me^{2\pi\xi}\ge1$ of the multiplicative coordinate; and $\widehat{F_{\rm raw}}$ is even, so
$\widehat{F_{\rm raw}}\ge0$ on $\RR$. Consequently
\[
\Xi(0)^2H_{\rm raw}(0)=\int_\RR\widehat{F_{\rm raw}}(\xi)\dd\xi=\frac1\pi\sum_{m\ge1}\frac{d(m)}{\sqrt m}\int_m^\infty\Gh_{\rm raw}(x)\,\frac{\dd x}x>0,
\]
so $H_{\rm raw}(0)>0$ and $C:=1/H_{\rm raw}(0)>0$.
\end{proposition}

\begin{proof}
By Proposition~\ref{prop:Hraw}(b),(c), $H_{\rm raw}\in\Wc_\delta$ for every $\delta\in(0,1)$, and its $\Gamma$-only transform is $\Gh_{\rm raw}$; fix
$\delta\in(0,\frac12)$. Lemma~\ref{lem:voronoi}(c) gives the series for every real $\xi$. In the multiplicative coordinate $\Gh_{\rm raw}\ge0$ on
$[1,\infty)$: for $x>1$ by \eqref{eq:Graw} and Theorem~\ref{thm:an-positive}, and $\Gh_{\rm raw}(1)=\frac1{32\pi^2}$. So every term is non-negative
when $\xi\ge0$. $F_{\rm raw}$ is even and real on $\RR$ (Proposition~\ref{prop:Hraw}(a)) and lies in $\Tc_\delta$ (Lemma~\ref{lem:voronoi}(a)); so
$\widehat{F_{\rm raw}}$ is real, even and integrable (Lemma~\ref{lem:I-decay}), and Fourier inversion holds at $t=0$, as $F_{\rm raw}$ and
$\widehat{F_{\rm raw}}$ are continuous and integrable. Since $\widehat{F_{\rm raw}}$ is continuous and non-negative, with
$\widehat{F_{\rm raw}}(0)=\int F_{\rm raw}=\frac1{32\pi^2}$ (Proposition~\ref{prop:C3}(b)), $\int\widehat{F_{\rm raw}}>0$. The series form follows from
Tonelli's theorem and the substitution $x=me^{2\pi\xi}$ in the $m$-th term. Finally $\Xi(0)$ is real, so $\Xi(0)^2H_{\rm raw}(0)>0$ forces
$H_{\rm raw}(0)>0$.
\end{proof}

Proposition~\ref{prop:posdef} is dual to Proposition~\ref{prop:C3}(b). The identity $\int F_{\rm raw}=\widehat{F_{\rm raw}}(0)=\frac1{32\pi^2}$ fixes
the sign of $C$ only once $H_{\rm raw}$ is known to have constant sign; the identity $F_{\rm raw}(0)=\int\widehat{F_{\rm raw}}$ needs only
$\widehat{F_{\rm raw}}\ge0$, which the positivity of the $a_n$ provides. For $F=CF_{\rm raw}$ the inequality $\hF\ge0$ is part of
Theorem~\ref{thm:main-S}(a); the point is that the proof above does not use the sign of $C$ (in particular it does not use
Theorem~\ref{thm:main-object}(b), which is stated for $H=CH_{\rm raw}$), so it determines that sign.

\begin{remark}[Checks of the normalising constant]\label{rem:sign}
The proofs of the main theorems take the sign of $C$ from Proposition~\ref{prop:posdef}. Two computations confirm it, Proposition~\ref{prop:C} and
the certified enclosure of $\Pi(0)>0$ in Theorem~\ref{thm:window} (as $H_{\rm raw}(0)=8\Pi(0)/\pi^2$ by Theorem~\ref{thm:F}); so does the identity of
Proposition~\ref{prop:C3}(b), given that $H_{\rm raw}$ has constant sign on $\RR$. The value of Proposition~\ref{prop:C} is
confirmed by the moment certificate of Section~\ref{ssec:moments} with the certified coefficients, which encloses $M_0=\int\widehat{F_{\rm raw}}$ and so
gives $H_{\rm raw}(0)=M_0/\Xi(0)^2=0.0011355178\pm1.5\cdot10^{-8}$; this enclosure contains that of Proposition~\ref{prop:C}. (A
floating-point evaluation of $M_0$, which is not part of the certificates, agrees with Proposition~\ref{prop:C} to $3.5\cdot10^{-12}$.)
\end{remark}

\subsection{Positivity near \texorpdfstring{$t=0$}{t = 0}}\label{ssec:moments}

Proposition~\ref{prop:posdef} is the case $t=0$ of a lower bound for $F_{\rm raw}(t)$ by the moments of $\widehat{F_{\rm raw}}$. For $j\ge0$ put
\begin{equation}\label{eq:moments}
M_{2j}:=\int_\RR\xi^{2j}\,\widehat{F_{\rm raw}}(\xi)\dd\xi=\frac1\pi\int_1^\infty W_j(x)\,\Gh_{\rm raw}(x)\,\frac{\dd x}x,
\end{equation}
where $W_j(x):=\sum_{m\le x}d(m)m^{-1/2}(\log(x/m)/2\pi)^{2j}$; the second form follows as the series form in Proposition~\ref{prop:posdef}. The moments are finite, since $\widehat{F_{\rm raw}}$ decays
exponentially (Lemma~\ref{lem:I-decay}), and positive.

\begin{lemma}[Moment minorant]\label{lem:moments}
For every real $t$ and every integer $K\ge0$,
\[
F_{\rm raw}(t)=\int_\RR\widehat{F_{\rm raw}}(\xi)\cos(2\pi\xi t)\dd\xi\ \ge\ P_K(t):=\sum_{j=0}^{2K+1}(-1)^j\frac{(2\pi t)^{2j}}{(2j)!}\,M_{2j}.
\]
If $P_K(t)>0$, then $\Xi(t)\ne0$ and $H_{\rm raw}(t)>0$.
\end{lemma}

\begin{proof}
The integral is Fourier inversion, as in the proof of Proposition~\ref{prop:posdef}, together with evenness. Let
$T_k(y):=\sum_{j\le k}(-1)^jy^{2j}/(2j)!$. Starting from $\cos y\le1$ and $\sin y\le y$ and integrating from $0$, induction on $k$ gives
$(-1)^{k+1}(\cos y-T_k(y))\ge0$ for $y\ge0$ and every $k\ge0$. For $k=2K+1$, where the truncation stops after a negative term, and by evenness,
$\cos y\ge T_{2K+1}(y)$ for every real $y$. Put $y=2\pi\xi t$ and integrate against $\widehat{F_{\rm raw}}\ge0$ (Proposition~\ref{prop:posdef}).
Finally $F_{\rm raw}(t)=\Xi(t)^2H_{\rm raw}(t)$ with $\Xi(t)$ real, so $F_{\rm raw}(t)>0$ forces $\Xi(t)\ne0$ and $H_{\rm raw}(t)>0$.
\end{proof}

\begin{theorem}[Moment certificate; computer-assisted]\label{thm:moments}
$P_7(t)>0$ for every real $t$ with $|t|\le10.355$; hence $H_{\rm raw}(t)>0$ there. The certificate uses the coefficients only through
$(a_n)\in\TB$: the lower bound for $P_7(t)$ that it certifies holds for every sequence in $\TB$ in place of $(a_n)$.
\end{theorem}

The certificate (Appendix~\ref{app:cert-M}) bounds $P_7(t)$ from below uniformly on $\TB$, through enclosures of the sixteen moments
$M_0,\dots,M_{30}$ of the centre sequence of $\TB$ and of the first eight coefficients separately, with all quadrature, truncation and tail
errors. One set of moments covers the whole interval, since $P_7$ is a polynomial in $t$; no Taylor model in $t$ is needed.
With the certified enclosures of the $a_n$ (Appendix~\ref{app:cert-coeffs}) in place of $\TB$, the same computation gives $P_7(t)>0$ for
$|t|\le10.402$. No $P_K$ can be positive at $t=14.1347\ldots$, the first zero of $\zeta$ on the critical line, where $F_{\rm raw}$ vanishes;
Theorem~\ref{thm:large-t-box} covers $|t|\ge9$.

\begin{remark}[Conditioning, and the elementary range]\label{rem:moments}
Lemma~\ref{lem:moments} is well conditioned near $t=0$. Every coefficient enters $M_0=F_{\rm raw}(0)$ with a positive weight, so the condition
number of $F_{\rm raw}(0)$ as a function of the coefficients is $1$, against about $6\cdot10^5$ for $\Pi(0)$ (Remark~\ref{rem:conditioning}); and
$\int\widehat{F_{\rm raw}}(\xi)|\cos(2\pi\xi t)|\dd\xi/F_{\rm raw}(t)$, which measures the cancellation in Lemma~\ref{lem:moments}, is $1.0000$,
$1.047$, $1.495$, $3.083$ and $9.34$ at $t=0,4,6,8,10$. This is why the box $\TB$, whose half-width for $a_1$ is $3\%$, suffices. With the
elementary bounds for $\Gh_{\rm raw}$ that Theorem~\ref{thm:an-bounds} and the bounds
$\sqrt{\pi/2y}\,e^{-y}(1-\frac1{8y})\le K_0(y)\le\sqrt{\pi/2y}\,e^{-y}$ give, instead of certified values of all the moments, $P_0$ shows
$H_{\rm raw}(t)>0$ for $|t|<4.4699$, and for $|t|<5.1995$ if $a_1$ is kept in its interval in $\TB$. These two ranges rest on two
one-dimensional integrals, a lower bound for $M_0$ and an upper bound for $M_2$, which we evaluated in Arb (Appendix~\ref{app:cert-M}); they are not
used in the proofs.
\end{remark}
